\documentclass[11pt,a4paper]{article}

\usepackage[utf8]{inputenc}
\usepackage[T1]{fontenc}
\usepackage[lithuanian]{babel}
\usepackage{lmodern}
\usepackage{geometry}
\usepackage{xcolor}
\usepackage{enumitem}
\usepackage{titlesec}
\usepackage{fancyhdr}
\usepackage{microtype}
\usepackage{ragged2e}

\geometry{
  top=1.8cm,
  bottom=1.8cm,
  left=2.1cm,
  right=2.1cm
}

\definecolor{titleblue}{HTML}{1F3A5F}
\definecolor{ruleblue}{HTML}{6EA8E5}
\definecolor{callout}{HTML}{EEF2F6}

\setlength{\parindent}{0pt}
\setlength{\parskip}{0.45em}
\renewcommand{\familydefault}{\sfdefault}

\titleformat{\section}
  {\large\bfseries\color{titleblue}}
  {}{0pt}{}
\titlespacing*{\section}{0pt}{0.9em}{0.25em}

\titleformat{\subsection}
  {\normalsize\bfseries\color{titleblue}}
  {}{0pt}{}
\titlespacing*{\subsection}{0pt}{0.65em}{0.2em}

\setlist[itemize]{leftmargin=1.3em,itemsep=0.05em,topsep=0.15em,parsep=0pt}

\pagestyle{fancy}
\fancyhf{}
\fancyfoot[C]{\thepage}
\renewcommand{\headrulewidth}{0pt}
\renewcommand{\footrulewidth}{0pt}

\newcommand{\calloutbox}[1]{%
  \begin{center}
  \colorbox{callout}{%
    \parbox{0.94\textwidth}{\centering\bfseries #1}%
  }
  \end{center}
}

\begin{document}

\begin{center}
  {\LARGE\bfseries\color{titleblue} Profesoriaus Jono Matulionio\\[2pt]
  jaunųjų matematikų konkursas}\par
  \vspace{0.35em}
  {\color{ruleblue}\rule{0.92\textwidth}{0.8pt}}\par
  \vspace{0.45em}
  {\itshape Trumpas gidas pradedančiajam olimpiadininkui}
\end{center}

\section{Kas tai per konkursas?}

\textbf{Profesoriaus Jono Matulionio jaunųjų matematikų konkursas} -- kasmet Kauno technologijos universitete (KTU) rengiamas respublikinis matematikos konkursas \textbf{9--12 klasių mokiniams}.

Tai vienas iš tradicinių Lietuvos matematikos konkursų. Jame susirenka keli šimtai matematikai gabių ir ja besidominčių mokinių iš visos Lietuvos. Uždavinius rengia KTU Matematikos ir gamtos mokslų fakulteto dėstytojai. Konkurso tikslas -- ne tik patikrinti mokyklines žinias, bet ir gebėjimą \textbf{mąstyti nestandartiškai, pastebėti idėją, ieškoti trumpesnio sprendimo ir argumentuoti savo samprotavimus}.

Matulionio konkursas ypač tinkamas mokiniui, kuris jau neblogai jaučiasi mokyklinėje matematikoje ir nori pradėti žengti į \textbf{olimpiadinę matematiką}.

\section{Kada ir kiek laiko vyksta konkursas?}

Konkursas tradiciškai vyksta \textbf{sausio pabaigoje KTU}. Uždaviniams spręsti skiriamos \textbf{3 valandos}.

Trys valandos atrodo daug, tačiau uždavinių taip pat nemažai. Todėl čia svarbus ne tik matematinis pajėgumas, bet ir gebėjimas \textbf{valdyti laiką}.

\section{Kokia konkurso struktūra?}

Konkurso užduotys sudarytos iš \textbf{dviejų dalių}.

\subsection{I dalis}

Pirmojoje dalyje pateikiama \textbf{10 trumpesnių uždavinių}.

Jie nėra paprasti mokykliniai pratimai. Dažnai reikia pastebėti gudrybę, dėsningumą ar netikėtą ryšį. Gali pasitaikyti algebros, geometrijos, skaičių teorijos, kombinatorikos ir loginio mąstymo uždavinių.

Šioje dalyje ypač svarbus \textbf{galutinis atsakymas}. Todėl nebūtina kiekvienam uždaviniui rašyti ilgo olimpiadinio įrodymo -- svarbiausia rasti teisingą rezultatą.

Tačiau I dalis turi dar vieną labai svarbią funkciją: ji yra tarsi \textbf{filtras į II dalies vertinimą}.

Kad būtų tikrinami II dalies sprendimai, I dalyje reikia surinkti \textbf{ne mažiau kaip 7 taškus iš 10}.

\calloutbox{Pirmasis tikslas: surinkti bent 7 taškus iš 10.}

Tačiau svarbu suprasti: \textbf{7 taškai dar nėra prizinės vietos rezultatas}. Tai pirmiausia reiškia, kad komisija pereis prie tavo II dalies sprendimų.

\subsection{II dalis}

Antrojoje dalyje pateikiami \textbf{2 sudėtingesni olimpiadinio pobūdžio uždaviniai}.

Čia jau nebeužtenka parašyti vien atsakymą.

Reikia parodyti:
\begin{itemize}
  \item kodėl tavo teiginiai teisingi;
  \item iš kur gaunamos formulės ar lygybės;
  \item kodėl nagrinėjami visi galimi atvejai;
  \item kaip iš vieno žingsnio logiškai pereinama prie kito.
\end{itemize}

Todėl II dalis tikrina jau kitą matematinį gebėjimą -- \textbf{mokėjimą sukurti ir užrašyti pilną matematinį argumentą}.

Net jei nepavyksta išspręsti viso uždavinio, verta užrašyti prasmingą pažangą: gerą lemos idėją, svarbią lygybę, įrodytą tarpinį teiginį ar teisingai išnagrinėtą atvejį.

\section{Kuo Matulionio konkursas ypatingas?}

Šis konkursas įdomus tuo, kad viename darbe susitinka \textbf{du skirtingi matematikos sprendimo būdai}.

I dalyje reikia būti greitam, lanksčiam ir mokėti pastebėti idėją.

II dalyje reikia sulėtėti, gilintis ir gebėti savo mintį paversti griežtu matematiniu įrodymu.

Todėl geras Matulionio konkurso dalyvis turi mokėti ir \textbf{greitai „medžioti“ trumpus uždavinius}, ir \textbf{ilgiau kovoti su vienu sudėtingu uždaviniu}.

Būtent dėl to šis konkursas yra gera tarpinė stotelė tarp įprastos mokyklinės matematikos ir rimtesnės olimpiadinės matematikos.

\section{Papildomas kelias į šalies etapą}

Matulionio konkursas svarbus dar ir tuo, kad \textbf{geriausiai jame pasirodę mokiniai gali gauti teisę dalyvauti Lietuvos mokinių matematikos olimpiados šalies etape}.

Tai reiškia, kad konkursas gali tapti \textbf{alternatyviu keliu į šalies etapą}: net jei miesto (savivaldybės) etape nepavyko pasirodyti pakankamai gerai, stiprus rezultatas Matulionio konkurse gali suteikti \textbf{kelialapį tiesiai į šalies etapą}, neperėjus per miesto etapą.

Todėl Matulionio konkursas gali būti ne tik gera treniruotė, bet ir labai svarbi galimybė visame olimpiadiniame sezone.

\section{Kokia turėtų būti strategija konkurso metu?}

Pradėjęs konkursą nebandyk būtinai spręsti uždavinių eilės tvarka.

Pirmiausia peržvelk I dalį ir ieškok uždavinių, kuriuose iš karto matai kryptį. Juos išspręsk pirmus.

Jeigu prie vieno uždavinio ilgai nejudi, palik jį ir eik toliau. I dalyje vienas sunkus uždavinys nėra vertas to, kad dėl jo prarastum galimybę išspręsti tris lengvesnius.

Pirmasis rimtas orientyras yra \textbf{7 taškai}. Pasiekęs šią ribą jau žinai, kad tavo II dalies darbas turės reikšmę.

Tačiau jei sieki gero rezultato, nereikėtų sustoti ties septyniais. Stiprus tikslas būtų \textbf{8--10 teisingų I dalies uždavinių}, o likusį laiką skirti II daliai.

II dalyje rinkis tą uždavinį, kuriame matai daugiau idėjų. Geriau vieną uždavinį išspręsti aiškiai ir beveik iki galo negu abiejuose palikti po kelias padrikas mintis.

Ir būtinai palik laiko patikrinti I dalies atsakymus. Viena aritmetinė klaida gali reikšti ne tik prarastą tašką -- jei esi ties riba, ji gali lemti, ar II dalies sprendimai apskritai bus tikrinami.

\section{Kaip ruoštis?}

Geriausias pasiruošimas -- \textbf{spręsti ankstesnių metų Matulionio konkurso uždavinius}.

Pradedančiajam nereikia iš karto bandyti išspręsti visko.

\calloutbox{10 I dalies uždavinių $\rightarrow$ 3 valandos $\rightarrow$ bandyti pasiekti 7/10.}

Kai 7/10 pradedi rinkti stabiliai, kelk kartelę iki \textbf{8/10}, paskui iki \textbf{9/10}.

Tuo pat metu mokykis spręsti II dalies uždavinius \textbf{be laiko spaudimo}. Iš pradžių vienam tokiam uždaviniui gali skirti net valandą ar kelias. Svarbiausia išmokti ne greitai gauti atsakymą, o surasti idėją ir ją tvarkingai įrodyti.

\section{Ką svarbiausia prisiminti pradedančiajam?}

Matulionio konkursas nėra egzaminas, kuriame privalai mokėti viską.

Jeigu pirmą kartą surinksi 4 ar 5 taškus iš 10, tai nėra nesėkmė. Tai parodo, \textbf{kur šiuo metu esi}.

Pirmasis aiškus tikslas -- pasiekti \textbf{7/10}.

Kitas -- stabiliai rinkti \textbf{8--9/10}.

O tada jau prasideda tikroji olimpiadinė užduotis: išmokti įveikti II dalies uždavinius.

Svarbiausia šiame etape ne medaliai. Svarbiausia išmokti tai, ko įprastoje matematikos pamokoje tenka gerokai mažiau: \textbf{ilgiau išbūti su nežinomu uždaviniu, bandyti skirtingas idėjas, klysti, grįžti atgal ir galiausiai pačiam atrasti sprendimą}.

\begin{center}
\textbf{Būtent nuo to ir prasideda olimpiadinė matematika.}
\end{center}

\end{document}
